% GENERATED FILE. Edit canonical lessons, then run npm run generate:books.
% canonical-source-hash: fnv1a64:57c526d64e3eff01
% canonical-lessons: ZH-C512-zhixiang, ZH-C512-zhuiqiu, ZH-C512-chengcai, ZH-C512-chengming, ZH-C512-mingxiao

\chapter{Dreams}
\label{ch:zh-dreams}
\hlchaptermodality{512}

\begin{chapteropening}
\textbf{By the end of this chapter:} \emph{I can talk about my ambition in life, chasing a dream, making something of myself, fame and top schools.}

Complete the last lesson of chapter 512: I can talk about my ambition in life, chasing a dream, making something of myself, fame and top schools.
\end{chapteropening}

\section[zhìxiàng]{zhìxiàng — ambition in life; the goal one sets oneself}

\label{lesson:ZH-C512-zhixiang}



\begin{quote}
Before the new one: say the Mandarin for one's own (by birth), then the Mandarin for personally (doing it oneself).
\end{quote}

\subsection*{You'll want to know: zhìxiàng}
\emph{zhìxiàng} — \textquotedblleft{}ambition in life; the goal one sets oneself\textquotedblright{}.

\begin{quote}
\emph{wǒ zuì dà de zhìxiàng shì dāng yīshēng} — My greatest ambition is to be a doctor.
\end{quote}

Its characters come later; for now it is a word for the ear.

\begin{cousinweb}
\emph{zhìxiàng} is built from \emph{zhì} 'a longing' and \emph{xiàng} 'to face', both inherited from Middle Chinese, \emph{tsyiH} and \emph{xjangH} (Guangyun, in Baxter's transcription), which glosses the first 'what the mind longs for' and the second 'to face'.
\end{cousinweb}

\begin{grammarlens}[title={What you've built}]
One more word for saying how it was.
\end{grammarlens}

\subsection*{Your turn}
\begin{itemize}
\raggedright
  \item \emph{Say it:} \emph{zhìxiàng}
  \item \emph{Say it:} \emph{zhìxiàng}, once more
  \item \emph{Say it:} \emph{wǒ zuì dà de zhìxiàng shì dāng yīshēng} — My greatest ambition is to be a doctor.
  \item \emph{Recall:} say the Mandarin for one's own (by birth), then the Mandarin for personally (doing it oneself), then say \emph{zhìxiàng} again
\end{itemize}

\begin{tcolorbox}[breakable,colback=teal!4,colframe=teal!35!black,title={Before you move on}]
What does zhìxiàng mean? (\textquotedblleft{}Ambition in life; the goal one sets oneself\textquotedblright{}.) Say it once more.
\end{tcolorbox}
\section[zhuīqiú]{zhuīqiú — to pursue (a goal); to chase a dream}

\label{lesson:ZH-C512-zhuiqiu}



\begin{quote}
Before the new one: say the Mandarin for personally (doing it oneself), then the Mandarin for one's ambition in life.
\end{quote}

\subsection*{You'll want to know: zhuīqiú}
\emph{zhuīqiú} — \textquotedblleft{}to pursue (a goal); to chase a dream\textquotedblright{}.

\begin{quote}
\emph{wǒ xiǎng zhuīqiú zìjǐ de mèngxiǎng} — I want to follow my own dream.
\end{quote}

Its characters come later; for now it is a word for the ear.

\begin{cousinweb}
\emph{zhuīqiú} is built from \emph{zhuī} 'to chase' and \emph{qiú} 'to seek', both inherited from Middle Chinese, \emph{trwij} and \emph{gjuw} (Guangyun, in Baxter's transcription), which glosses the first 'to chase' and the second 'to search for'.
\end{cousinweb}

\begin{grammarlens}[title={What you've built}]
One more word for saying how it was.
\end{grammarlens}

\subsection*{Your turn}
\begin{itemize}
\raggedright
  \item \emph{Say it:} \emph{zhuīqiú}
  \item \emph{Say it:} \emph{zhuīqiú}, once more
  \item \emph{Say it:} \emph{wǒ xiǎng zhuīqiú zìjǐ de mèngxiǎng} — I want to follow my own dream.
  \item \emph{Recall:} say the Mandarin for personally (doing it oneself), then the Mandarin for one's ambition in life, then say \emph{zhuīqiú} again
\end{itemize}

\begin{tcolorbox}[breakable,colback=teal!4,colframe=teal!35!black,title={Before you move on}]
What does zhuīqiú mean? (\textquotedblleft{}To pursue (a goal); to chase a dream\textquotedblright{}.) Say it once more.
\end{tcolorbox}
\section[\texorpdfstring{\zh{成才} (chéngcái)}{成才 (chéngcái)}]{\zh{成才} (chéngcái) — to make something of oneself}

\label{lesson:ZH-C512-chengcai}



\begin{quote}
Before the new one: say the Mandarin for one's ambition in life, then the Mandarin for to pursue (a goal).
\end{quote}

\subsection*{You'll want to know: \zh{成才}}
\textbf{\zh{成才}} — \emph{chéngcái} — \textquotedblleft{}to make something of oneself\textquotedblright{}.

\begin{quote}
\emph{tā cóngxiǎo nǔlì xuéxí hòulái chéngcái le} — He studied hard from childhood and later made something of himself.
\end{quote}

\begin{cousinweb}
\textbf{\zh{成才}} (\emph{chéngcái}) is built from \zh{成} 'to become' and \zh{才} 'talent', both inherited from Middle Chinese, \emph{dzyeng} and \emph{dzoj} (Guangyun, in Baxter's transcription), which glosses the first 'to complete' and the second 'ability' and 'a gift for letters'.
\end{cousinweb}

\begin{grammarlens}[title={What you've built}]
One more word for saying how it was.
\end{grammarlens}

\subsection*{Your turn}
\begin{itemize}
\raggedright
  \item \emph{Say it:} \emph{chéngcái}
  \item \emph{Say it:} \emph{chéngcái}, once more
  \item \emph{Say it:} \emph{tā cóngxiǎo nǔlì xuéxí hòulái chéngcái le} — He studied hard from childhood and later made something of himself.
  \item \emph{Recall:} say the Mandarin for one's ambition in life, then the Mandarin for to pursue (a goal), then say \emph{chéngcái} again
\end{itemize}

\begin{tcolorbox}[breakable,colback=teal!4,colframe=teal!35!black,title={Before you move on}]
What does \zh{成才} mean? (\textquotedblleft{}To make something of oneself\textquotedblright{}.) Say it once more.
\end{tcolorbox}
\section[\texorpdfstring{\zh{成名} (chéngmíng)}{成名 (chéngmíng)}]{\zh{成名} (chéngmíng) — to become famous}

\label{lesson:ZH-C512-chengming}



\begin{quote}
Before the new one: say the Mandarin for to pursue (a goal), then the Mandarin for to make something of oneself.
\end{quote}

\subsection*{You'll want to know: \zh{成名}}
\textbf{\zh{成名}} — \emph{chéngmíng} — \textquotedblleft{}to become famous\textquotedblright{}.

\begin{quote}
\textbf{\zh{成名以后他感觉很不开心}} \emph{chéngmíng yǐhòu tā gǎnjué hěn bù kāixīn} — After he became famous, he felt very unhappy.
\end{quote}

\begin{cousinweb}
\textbf{\zh{成名}} (\emph{chéngmíng}) is built from \zh{成} 'to complete' and \zh{名} 'a name', both inherited from Middle Chinese, \emph{dzyeng} and \emph{mjieng} (Guangyun, in Baxter's transcription), which glosses the second 'a name' and 'a title'.
\end{cousinweb}

\begin{grammarlens}[title={What you've built}]
One more word for saying how it was.
\end{grammarlens}

\subsection*{Your turn}
\begin{itemize}
\raggedright
  \item \emph{Say it:} \emph{chéngmíng}
  \item \emph{Say it:} \emph{chéngmíng}, once more
  \item \emph{Say it:} \emph{chéngmíng yǐhòu tā gǎnjué hěn bù kāixīn} — After he became famous, he felt very unhappy.
  \item \emph{Recall:} say the Mandarin for to pursue (a goal), then the Mandarin for to make something of oneself, then say \emph{chéngmíng} again
\end{itemize}

\begin{tcolorbox}[breakable,colback=teal!4,colframe=teal!35!black,title={Before you move on}]
What does \zh{成名} mean? (\textquotedblleft{}To become famous\textquotedblright{}.) Say it once more.
\end{tcolorbox}
\section[\texorpdfstring{\zh{名校} (míngxiào)}{名校 (míngxiào)}]{\zh{名校} (míngxiào) — top school; famous university}

\label{lesson:ZH-C512-mingxiao}



\begin{quote}
Before the new one: say the Mandarin for to make something of oneself, then the Mandarin for to become famous.
\end{quote}

\subsection*{You'll want to know: \zh{名校}}
\textbf{\zh{名校}} — \emph{míngxiào} — \textquotedblleft{}top school; famous university\textquotedblright{}.

\begin{quote}
\textbf{\zh{我想考上名校}} \emph{wǒ xiǎng kǎoshàng míngxiào} — I want to get into a top school.
\end{quote}

\begin{cousinweb}
\textbf{\zh{名校}} (\emph{míngxiào}) is built from \zh{名} 'a name' and \zh{校} 'a school', both inherited from Middle Chinese, \emph{mjieng} and \emph{haewH} (Guangyun, in Baxter's transcription), which glosses the second as the title of an officer; the sense 'a school' is not given there.
\end{cousinweb}

\begin{grammarlens}[title={What you've built}]
One more word for saying how it was.
\end{grammarlens}

\subsection*{Your turn}
\begin{itemize}
\raggedright
  \item \emph{Say it:} \emph{míngxiào}
  \item \emph{Say it:} \emph{míngxiào}, once more
  \item \emph{Say it:} \emph{wǒ xiǎng kǎoshàng míngxiào} — I want to get into a top school.
  \item \emph{Recall:} say the Mandarin for to make something of oneself, then the Mandarin for to become famous, then say \emph{míngxiào} again
\end{itemize}

\begin{tcolorbox}[breakable,colback=teal!4,colframe=teal!35!black,title={Before you move on}]
What does \zh{名校} mean? (\textquotedblleft{}Top school; famous university\textquotedblright{}.) Say it once more.
\end{tcolorbox}
